\chapter{Решение четырёх блоков}
\label{ch:practice}

Методичка делит работу на два блока (HTML и CSS), внутри каждого --- два
подблока. Итого четыре независимых задания, каждое --- отдельный pen на
Codepen. Соответствие подблоков и pen-ов:

\begin{center}
\begin{tabular}{|c|l|l|l|}
\hline
\textbf{Подблок} & \textbf{Тема} & \textbf{Pen (slug)} & \textbf{Short URL} \\
\hline
1.1 & Теги для разметки текста и атрибуты            & \texttt{wvVrwjg} & \texttt{clck.ru/3E73mG} \\
1.2 & Списки и таблицы                                & \texttt{KKOXPBZ} & \texttt{clck.ru/3E73xM} \\
2.1 & Селекторы и свойства                            & \texttt{MWNEgBG} & \texttt{clck.ru/3E744t} \\
2.2 & Оформление текстовых блоков при помощи CSS      & \texttt{eYqGOjj} & \texttt{clck.ru/3E74AH} \\
\hline
\end{tabular}
\end{center}

\section{Блок 1.1: Теги для разметки текста}

Исходный pen содержит три секции, в каждой требуется выполнить отдельную
задачу разметки.

\subsection{Задача 1 --- статья про нейросеть}

Текст про нейросеть в первом \texttt{<section>} был представлен без
структуры (просто два абзаца текста, разделённые переводами строк).
Требовалось:
\begin{enumerate}
  \item Превратить предложение \textit{<<Что такое нейросеть>>} в
        заголовок второго уровня (\texttt{<h2>}).
  \item Выделить два следующих абзаца тегами \texttt{<p>}.
  \item Во втором абзаце слово \textit{<<научилась>>} превратить
        в гиперссылку на статью о работе нейросети Яндекса со
        стихами Егора Летова.
\end{enumerate}

Готовая разметка:
\begin{lstlisting}[language=HTML]
<section>
  <h2>Что такое нейросеть</h2>
  <p>Грубо говоря, это сложная программа, состоящая из многих
     программ поменьше, ...</p>
  <p>Недавно, например, нейросеть
     <a href="https://www.sostav.ru/publication/nejroset-yandeksa
     -napisala-tekst-dlya-muzykalnogo-alboma-22860.html">научилась</a>
     писать песни в стиле Егора Летова. ...</p>
</section>
\end{lstlisting}

\subsection{Задача 2 --- разметка <<Дорога>>}

В исходном pen-е разметка второго \texttt{<section>} содержала
четыре ошибки: закрывающий \texttt{</h2>} без открывающего тега для
слова \textit{<<Дорога>>}; \texttt{<blockquote>} с двумя
закрывающими тегами \texttt{</blockquote>} и \texttt{</p>} в
неправильном порядке. Требовалось исправить разметку так, чтобы
получилась последовательность: заголовок --- абзац --- цитата ---
абзац.

Готовая разметка:
\begin{lstlisting}[language=HTML]
<section>
  <h2>Дорога</h2>
  <p>Вечером второго дня я свернул в деревню Мшага ...</p>
  <blockquote>
    <p>Вдоль трассы М10 очень много мёртвых или полумёртвых
       деревень ...</p>
  </blockquote>
  <p>После того как я проехал Великий Новгород, дожди
     закончились ...</p>
</section>
\end{lstlisting}

\subsection{Задача 3 --- гейзеры Камчатки}

Третий \texttt{<section>} содержал заголовок второго уровня
\textit{<<Источники и гейзеры Камчатки>>} (уже корректно
размеченный) и три коротких рассказа --- про посёлок Малки,
Налычевскую долину и Паратунку. Требовалось:
\begin{itemize}
  \item Превратить названия гейзеров в заголовки третьего уровня
        (\texttt{<h3>}).
  \item Описание каждого гейзера обернуть в \texttt{<p>}.
  \item Перед каждым описанием вставить тег \texttt{<img>} с
        соответствующей фотографией и атрибутом \texttt{alt},
        равным названию гейзера.
\end{itemize}

Файлы изображений лежат в \texttt{lab12/codepen-solutions/img/} и
подключаются относительными ссылками:
\begin{lstlisting}[language=HTML]
<h3>Посёлок Малки</h3>
<img src="../img/malki.jpg" alt="Посёлок Малки">
<p>Малки – посёлок в 132 км от Петропавловска-Камчатского. ...</p>

<h3>Налычевская долина</h3>
<img src="../img/nalychevo.jpg" alt="Налычевская долина">
<p>...</p>

<h3>Паратунка</h3>
<img src="../img/paratunka.jpg" alt="Паратунка">
<p>...</p>
\end{lstlisting}

% — Скриншот: рендер решения блока 1.1 —
\begin{figure}[h]
  \centering
  \includegraphics[width=0.85\textwidth]{img/01_block_1_1}
  \caption{Рендер решения блока 1.1: разметка текста, ссылка, фотографии гейзеров}
  \label{fig:block_1_1}
\end{figure}

Готовый форк pen-а на Codepen:
\begin{lstlisting}
https://codepen.io/blxckweed/pen/EaNmYge
\end{lstlisting}

\section{Блок 1.2: Списки и таблицы}

Исходный pen содержит четыре секции, в каждой --- отдельная задача на
списки или таблицы.

\subsection{Задача 1 --- маркированный список}

В первой секции были перечислены пять имён через запятую:
Мартин Лютер Кинг, Роберт Оппенгеймер, Альфред Хичкок, Чарли Чаплин,
Коко Шанель. Требовалось превратить перечисление в маркированный
список \texttt{<ul>}.

\begin{lstlisting}[language=HTML]
<ul>
  <li>Мартин Лютер Кинг</li>
  <li>Роберт Оппенгеймер</li>
  <li>Альфред Хичкок</li>
  <li>Чарли Чаплин</li>
  <li>Коко Шанель</li>
</ul>
\end{lstlisting}

\subsection{Задача 2 --- упорядоченный список миллиардеров}

ТОП-5 списка Forbes требовалось расположить в порядке \emph{убывания}
состояния. Корректный порядок: Гейтс (\$75~млрд), Ортега
(\$67~млрд), Баффет (\$60{,}8~млрд), Слим (\$50~млрд), Безос
(\$45{,}2~млрд). Поскольку порядок важен, использован
\texttt{<ol>}, а не \texttt{<ul>}.

\begin{lstlisting}[language=HTML]
<ol>
  <li>Билл Гейтс ($75 млрд)</li>
  <li>Амансио Ортега ($67 млрд)</li>
  <li>Уоррен Баффет ($60,8 млрд)</li>
  <li>Карлос Слим Элу ($50 млрд)</li>
  <li>Джефф Безос ($45,2 млрд)</li>
</ol>
\end{lstlisting}

\subsection{Задача 3 --- таблица заказа аниматоров}

Дано семь строк <<персонаж/время/цена>>. Требовалось построить
таблицу с шапкой <<Персонаж~/~Время выступления~/~Цена>> и итоговой
строкой <<ИТОГО>>, занимающей две первые ячейки, с общей суммой в
третьей. Сумма проверяется арифметически:
$1500+4800+8400+5200+3300+4000+2150=29\,350$~руб.

Также требовалось использовать <<все теги, изученные на занятии>>:
\texttt{<table>}, \texttt{<thead>}, \texttt{<tbody>},
\texttt{<tfoot>}, \texttt{<tr>}, \texttt{<th>}, \texttt{<td>},
\texttt{colspan}.

\begin{lstlisting}[language=HTML]
<table>
  <thead>
    <tr>
      <th>Персонаж</th><th>Время выступления</th><th>Цена</th>
    </tr>
  </thead>
  <tbody>
    <tr><td>Белоснежка</td><td>12:00</td><td>1 500 руб.</td></tr>
    <tr><td>Лунтик</td><td>12:30</td><td>4 800 руб.</td></tr>
    <tr><td>Маша и Медведь</td><td>13:15</td><td>8 400 руб.</td></tr>
    <tr><td>Бэтмен</td><td>13:45</td><td>5 200 руб.</td></tr>
    <tr><td>Джек Воробей</td><td>14:30</td><td>3 300 руб.</td></tr>
    <tr><td>Аватар</td><td>14:50</td><td>4 000 руб.</td></tr>
    <tr><td>Энгри Бердс</td><td>15:20</td><td>2 150 руб.</td></tr>
  </tbody>
  <tfoot>
    <tr>
      <td colspan="2">ИТОГО</td>
      <td>29 350 руб.</td>
    </tr>
  </tfoot>
</table>
\end{lstlisting}

\subsection{Задача 4 --- починка таблицы звёзд эстрады}

В исходной таблице были три типа ошибок:
\begin{enumerate}
  \item Лишние пустые ячейки во второй строке шапки --- из-за
        \texttt{rowspan="2"} в первой строке колонки <<Возраст>>
        и <<Знак зодиака>> уже занимают вторую строку, и
        пустые \texttt{<td>} там избыточны.
  \item Лишний пустой \texttt{<td>} в строке Челентано.
  \item Атрибут \texttt{rowspan="2"} у ячейки <<Близнецы>> в
        строке Азнавура без соответствующего сокращения числа
        ячеек в строке Маккартни --- из-за этого следующая
        строка получает лишний \texttt{<td>}, который смещает
        остальные данные.
  \item Лишний пустой \texttt{<td>} в строке Мадонны.
\end{enumerate}

Исправленная таблица --- ровно семь строк и четыре колонки
(2 строки шапки + 5 строк данных):

\begin{lstlisting}[language=HTML]
<table>
  <thead>
    <tr>
      <th colspan="2">ФИО</th>
      <th rowspan="2">Возраст</th>
      <th rowspan="2">Знак зодиака</th>
    </tr>
    <tr>
      <th>Имя</th>
      <th>Фамилия</th>
    </tr>
  </thead>
  <tbody>
    <tr><td>Майкл</td><td>Джексон</td><td>51</td><td>Дева</td></tr>
    <tr><td>Адриано</td><td>Челентано</td><td>78</td><td>Козерог</td></tr>
    <tr><td>Шарль</td><td>Азнавур</td><td>94</td><td>Близнецы</td></tr>
    <tr><td>Пол</td><td>Маккартни</td><td>74</td><td>Близнецы</td></tr>
    <tr><td colspan="2">Мадонна</td><td>58</td><td>Лев</td></tr>
  </tbody>
</table>
\end{lstlisting}

% — Скриншот: рендер решения блока 1.2 —
\begin{figure}[h]
  \centering
  \includegraphics[width=0.85\textwidth]{img/02_block_1_2}
  \caption{Рендер решения блока 1.2: два списка и две таблицы}
  \label{fig:block_1_2}
\end{figure}

Готовый форк pen-а на Codepen:
\begin{lstlisting}
https://codepen.io/blxckweed/pen/YPpVKpm
\end{lstlisting}

\section{Блок 2.1: Селекторы и свойства}

Исходный pen содержит размеченный текст известной микро-сказки
Л.\,Петрушевской <<Пуськи Бятые>>: \textit{<<Сяпала Калуша с
Калушатами по напушке~...>>}. Разметку трогать запрещено --- задача
исключительно в написании четырёх CSS-правил.

\subsection{Задача 1 --- синий заголовок <<Калуша с Калушатами>>}

Главная сложность: на странице два \texttt{<h2>} --- один в шапке
секции (<<Разметка к задачам:>>) и один вложенный в
\texttt{<div>} (<<Калуша с Калушатами>>). Простой селектор
\texttt{h2} перекрасил бы оба. Нужен селектор, попадающий только во
вложенный заголовок:

\begin{lstlisting}[language=CSS]
section > div > h2 {
  color: blue;
}
\end{lstlisting}

\subsection{Задача 2 --- коричневые упоминания Калушат}

Все упоминания Калушат в тексте обёрнуты в \texttt{<strong>}.
Поскольку других \texttt{<strong>} в pen-е нет, достаточно
селектора по тегу:

\begin{lstlisting}[language=CSS]
strong {
  color: brown;
}
\end{lstlisting}

\subsection{Задача 3 --- жёлтая жирная Бутявка в цитатах}

Слово <<Бутявка>> (и его словоформы) встречается в тексте многократно
и обёрнуто в \texttt{<span>}. Нужно стилизовать только те
\texttt{<span>}, что находятся внутри цитат \texttt{<blockquote>} ---
без перекраски Бутявок в обычных абзацах:

\begin{lstlisting}[language=CSS]
blockquote span {
  color: yellow;
  font-weight: bold;
}
\end{lstlisting}

\subsection{Задача 4 --- курсивные Калушата в списке}

В нумерованном списке \texttt{<ol>} тоже встречаются Калушата ---
к ним нужно применить курсивное начертание дополнительно к уже
заданному коричневому цвету:

\begin{lstlisting}[language=CSS]
ol strong {
  font-style: italic;
}
\end{lstlisting}

% — Скриншот: рендер решения блока 2.1 —
\begin{figure}[h]
  \centering
  \includegraphics[width=0.85\textwidth]{img/03_block_2_1}
  \caption{Рендер решения блока 2.1: разноцветный текст, жёлтые Бутявки в цитатах, курсив Калушат в списке}
  \label{fig:block_2_1}
\end{figure}

Готовый форк pen-а на Codepen:
\begin{lstlisting}
https://codepen.io/blxckweed/pen/vEymBmY
\end{lstlisting}

\section{Блок 2.2: Оформление текстовых блоков}

Исходный pen содержит размеченную энциклопедическую статью про
А.~Эйнштейна (заголовок, портрет, лид-абзац, несколько обычных
абзацев, список научных работ, цитата). Задача --- написать пять
CSS-правил.

\subsection{Задача 1 --- главный заголовок}

Заголовок с классом \texttt{main-header}: размер 56\,px, цвет
\texttt{\#365f91} (тёмно-синий), шрифт с засечками.

\begin{lstlisting}[language=CSS]
.main-header {
  font-size: 56px;
  color: #365f91;
  font-family: serif;
}
\end{lstlisting}

\subsection{Задача 2 --- размер всех абзацев}

\begin{lstlisting}[language=CSS]
p {
  font-size: 20px;
}
\end{lstlisting}

\subsection{Задача 3 --- подчёркивание в лид-абзаце}

Внутри абзаца с классом \texttt{lead} (первый абзац --- расширенное
определение) слова \textit{<<Альбе́рт Эйнште́йн>>} выделены тегом
\texttt{<strong>}. Нужно добавить им подчёркивание:

\begin{lstlisting}[language=CSS]
p.lead strong {
  text-decoration: underline;
}
\end{lstlisting}

\subsection{Задача 4 --- два уровня маркеров списка}

Список научных работ Эйнштейна --- \texttt{<ul~class="works-list">}.
Внутри первого \texttt{<li>} есть вложенный \texttt{<ul>} с одним
пунктом про закон взаимосвязи массы и энергии. Требуется:
\begin{itemize}
  \item все пункты обоих уровней --- цвета \texttt{\#365f91};
  \item первый уровень --- маркер <<квадрат>> ($\blacksquare$);
  \item второй уровень --- маркер <<пустой круг>> ($\bigcirc$).
\end{itemize}

\begin{lstlisting}[language=CSS]
.works-list {
  color: #365f91;
  list-style-type: square;
}

.works-list ul {
  list-style-type: circle;
}
\end{lstlisting}

\subsection{Задача 5 --- цитата с фоновой картинкой}

Цитата \texttt{<blockquote>} в конце статьи --- про вклад
Эйнштейна в популяризацию физики. Требовалось:
\begin{itemize}
  \item установить фоновую картинку (по ссылке методички ---
        изображение сетки/паттерна);
  \item изменить цвет текста на \texttt{\#494429};
  \item включить курсивное начертание и шрифт с засечками.
\end{itemize}

Картинка фона лежит локально в \texttt{codepen-solutions/img/blockquote\_bg.jpg}
(скачана с Яндекс.Диска по ссылке методички); в Codepen её можно
загрузить как \texttt{Asset} через
\texttt{Assets~\(\rightarrow\)~Image} и подставить полученный URL.

\begin{lstlisting}[language=CSS]
blockquote {
  background-image: url("../img/blockquote_bg.jpg");
  color: #494429;
  font-style: italic;
  font-family: serif;
}
\end{lstlisting}

% — Скриншот: рендер решения блока 2.2 —
\begin{figure}[h]
  \centering
  \includegraphics[width=0.85\textwidth]{img/04_block_2_2}
  \caption{Рендер решения блока 2.2: крупный заголовок, подчёркнутый лид, цветной список с двумя уровнями маркеров, цитата на фоновой картинке}
  \label{fig:block_2_2}
\end{figure}

Готовый форк pen-а на Codepen:
\begin{lstlisting}
https://codepen.io/blxckweed/pen/vEymBxB
\end{lstlisting}

\section{Готовые форки в Codepen}

Сводная таблица ссылок на форки, сохранённые под аккаунтом студента:

\begin{center}
\begin{tabular}{|c|l|l|}
\hline
\textbf{Блок} & \textbf{Тема} & \textbf{URL форка} \\
\hline
1.1 & Теги для разметки текста        & \url{https://codepen.io/blxckweed/pen/EaNmYge} \\
1.2 & Списки и таблицы                & \url{https://codepen.io/blxckweed/pen/YPpVKpm} \\
2.1 & Селекторы и свойства            & \url{https://codepen.io/blxckweed/pen/vEymBmY} \\
2.2 & Оформление текстовых блоков     & \url{https://codepen.io/blxckweed/pen/vEymBxB} \\
\hline
\end{tabular}
\end{center}

\textbf{Самопроверка:}
\begin{itemize}
  \item все четыре стартовых pen-а форкнуты под собственным
        аккаунтом Codepen, изменения сохранены;
  \item четыре URL форков подставлены в таблицу выше и в листинги
        внутри глав;
  \item локальные \texttt{index.html} в \texttt{lab12/codepen-solutions/}
        открываются в браузере и визуально совпадают с
        Codepen-форками;
  \item все четыре скриншота (\texttt{01\_block\_1\_1.png} ---
        \texttt{04\_block\_2\_2.png}) положены в
        \texttt{latex-report/img/}, имена совпадают с
        \texttt{\textbackslash includegraphics} в этой главе.
\end{itemize}

\endinput
